 \setlength{\baselineskip}{20pt}
\chapter{绪论}
\label{cha:intro}

\section{研究背景及意义}

\subsection{铁路道岔伤损检测的工程需求}

铁路道岔是列车实现转线、越行和折返等运行组织功能的关键基础设施，由尖轨、基本轨、翼轨、护轨、转辙机构及相关连接部件组成。与普通直线钢轨相比，道岔结构更加复杂，局部几何突变明显，轮轨接触关系和受力状态更为恶劣，是铁路线路中典型的结构薄弱环节。随着我国高速铁路和既有铁路网络规模持续扩大，道岔设备的服役数量、服役强度和维护压力均显著增加，其健康状态直接关系到列车运行安全和线路运输效率。

截至 2025 年底，我国铁路营业里程已达到 16.5 万公里，其中高速铁路营业里程达到 5.04 万公里，稳居世界首位\upcite{xinhuanet2026railway}。在高速、大密度运行条件下，道岔尖轨、基本轨及辙叉等关键部件长期承受轮轨冲击、交变弯曲、振动磨耗和转辙动作反复加载，容易产生疲劳裂纹、剥离掉块、压溃磨耗等表面及近表面伤损。特别是尖轨尖端及密贴区域，由于截面薄、受力集中、安装空间狭小，伤损早期信号微弱且不易观测，一旦未能及时发现并处置，可能进一步发展为断轨、挤岔、脱轨等严重安全事故。

\begin{figure}[htbp]
    \centering
    \includegraphics[width=0.64\textwidth]{figures/turnout_structure.jpg}
    \figcaption{铁路道岔结构图}{Schematic diagram of the railway turnout structure}
    \label{fig:turnout_structure}
\end{figure}

图 \ref{fig:turnout_structure} 给出了铁路道岔的典型结构图。可以看出，道岔区域内部构件密集，空间约束明显，且不同部件之间存在复杂的机械连接关系。这使得常规检测设备，尤其是依赖接触式或需较大安装空间的传感装置，难以在道岔关键薄弱区域进行长期固定布设，且在狭窄空间与转辙动作干扰下极易产生漏检或误报，从而增加了现场巡检和状态监测的难度。

围绕钢轨和道岔伤损检测，工程领域已经发展出超声、声发射、机器视觉、漏磁、涡流以及交流电磁场检测等多种无损检测技术\upcite{xiong2023detection}。这些方法在普通钢轨探伤、焊缝检测和表面缺陷识别等方面已取得较多应用，但当其应用于道岔尖轨、密贴区及转辙结构附近时，仍存在明显局限。
 
超声检测技术具有较强的内部缺陷探测能力，在钢轨探伤中应用广泛。然而，传统压电超声通常需要耦合剂，探头尺寸和安装姿态对检测结果影响较大，难以适应道岔尖端狭窄区域的长期固定布设。超声导波能够实现较长距离传播，但在道岔复杂结构中会受到滑床板、连接件和截面突变等因素影响，易产生多模态散射和波形混叠，导致缺陷定位和定量困难\upcite{chengchao2025ultrasound}。
 
声发射技术通过监听裂纹扩展过程中释放的弹性波实现损伤监测，属于被动式检测方法。其优点是能够捕捉裂纹扩展等动态过程，但在道岔现场，转辙动作、列车通过、轮轨冲击和周边机械振动都会产生强噪声干扰，容易造成误报或漏报\upcite{WEI2023SHANGJIAN}。
 
机器视觉方法适合进行表面可见缺陷识别，具有非接触、直观和易于部署等优势。然而，视觉检测依赖清晰光学视场，容易受到油污、雨雪、粉尘和光照变化影响，并且无法直接感知金属内部或近表面的早期裂纹\upcite{TDYY202504010}。
 
漏磁、涡流和交流电磁场检测等电磁无损检测技术对金属表面及近表面缺陷较为敏感，在钢轨和金属构件检测中具有较好应用基础。但是，传统磁学检测设备往往依赖体积较大的励磁结构或扫描式探头，难以在道岔狭窄区域实现长期、稳定、低扰动布设；同时，铁路现场存在转辙机、信号设备和牵引供电系统等复杂电磁干扰源，在强背景场下提取由早期微小缺陷引起的微弱电磁异常信号变得极具挑战性；此外，单一电磁场测量往往难以克服软场效应与逆问题多解性，导致缺陷定位困难\upcite{jiangjiahong2025electromagnetic,wangkang2025based}。
 
综上，现有钢轨探伤技术虽然各有优势，但在道岔伤损检测场景下，普遍面临安装空间受限、抗干扰能力不足、在线监测困难或早期近表面缺陷识别能力有限等问题。因此，有必要探索一种更适合道岔复杂结构和现场环境的新型无损检测方案。

\subsection{电磁层析成像用于道岔伤损检测的可行性与意义}

电磁层析成像（Electromagnetic Tomography, EMT）是一种基于交变电磁场感应效应的非接触式无损检测技术。其基本思想是通过外部激励线圈在被测区域建立交变电磁场，并利用传感器阵列测量目标内部介质分布变化所引起的边界电磁响应扰动，进而反演被测区域内部电导率、磁导率或与缺陷相关的等效参数分布。
 
与超声、视觉等检测方式相比，EMT 在道岔伤损检测中具有以下潜在优势。首先，EMT 传感器可设计为平面线圈阵列或柔性共形结构，能够较好适应道岔表面及尖轨邻近区域的狭小空间，为贴附式、便携式或原位式检测提供结构基础\upcite{TSE201948}。其次，EMT 通过测量导电导磁材料中交变电磁场的扰动响应，理论上能够对金属表层及近表层伤损形成敏感观测，适合用于疲劳裂纹、剥离、局部缺口等缺陷检测。再次，EMT 不依赖机械耦合剂和光学视场，对机械振动和表面污染的敏感性相对较低，更适合铁路现场复杂环境。
 
然而，将 EMT 应用于道岔伤损检测仍面临两方面关键难题。第一，高导电率、高磁导率铁磁材料背景下，强背景场掩盖了伤损引起的微弱电磁扰动，亟需高稳定度激励、高信噪比模拟前端和可靠的复数响应采集系统以提升信噪比；第二，EMT 逆问题具有显著病态性和非线性，特别是在强趋肤效应与磁电耦合作用下，传统基于灵敏度矩阵的线性重建方法极易产生伪影且泛化能力不足，难以在复杂铁磁场景中获得高分辨率的检测结果。第三，道岔伤损检测的工程应用场景要求系统具备的小型化与便携性，这对 EMT 系统硬件集成设计提出了更高要求。
 
因此，围绕道岔伤损 EMT 检测开展研究，需要同时解决硬件平台与反演算法两个层面的问题：在硬件层面，需要构建便携式、高信噪比、多通道复用的 EMT 数据采集平台；在算法层面，需要发展能够充分利用多视角复数响应、并适应铁磁材料强非线性特征的智能反演模型。本文正是在这一工程背景与理论需求下展开研究。

\section{国内外研究现状}

	\subsection{电磁层析成像系统研究概况}
 
EMT 技术自 20 世纪 90 年代提出以来，已逐步从医学成像、工业过程监测扩展到金属无损检测、结构健康监测和复杂介质参数反演等方向\upcite{WANG2020,9405663,9109633}。从发展脉络看，EMT 系统研究主要围绕传感器阵列拓扑、激励与测量硬件、信号调理链路、图像重建算法以及应用场景适配展开。
 
国外相关研究起步较早。1993 年，S. Al-Zeibak 等人构建了用于生物医学成像的 EMT 实验系统，验证了利用磁场响应反演电导率分布的可行性；同年，UMIST 的 A. J. Peyton 团队面向工业多相流检测提出了基于平行场激励的 EMT 系统原型，为工业 EMT 系统发展奠定了基础\upcite{al1993feasibility,zzyu1993emt}。随后，Peyton 团队提出的 16 线圈分时复用结构成为 EMT 阵列设计中的经典方案，对后续多通道 EMT 系统具有重要参考意义\upcite{peyton1995electromagnetic}。
 
图 \ref{fig:peyton_emt_system} 给出了 Peyton 团队早期环形线圈分时复用 EMT 系统示意。该系统通过线圈轮流激励与接收，获得多视角边界响应，为后续阵列式 EMT 系统提供了典型结构范式。

\begin{figure}[htbp]
    \centering
    \includegraphics[width=0.8\textwidth]{figures/emt1995.jpg}
    \figcaption{Peyton 的环形线圈分时复用 EMT 系统}{Peyton's ring-shaped time-division multiplexed EMT system}
    \label{fig:peyton_emt_system}
\end{figure}

随着应用对象不断拓展，国外学者进一步提出了 H 型阵列、共面阵列、C 型阵列、柔性阵列和无线供能式传感结构等多种专用传感器拓扑，以适应薄板、轧钢、复杂异形金属件以及封闭空间等不同检测对象\upcite{yin2006planar,ma2008development,daura2020wireless,fauchard2024electromagnetic}。近年来，国外研究重点逐渐由“硬件可实现性”转向“高精度反演与智能解释”，多频 EMT、磁感应层析成像与深度学习模型的结合成为重要趋势。
 
国内 EMT 研究虽起步稍晚，但在传感器理论、系统数字化实现和轨道交通专用化应用方面形成了较为系统的发展。天津大学等单位较早开展了传感器参数优化、近似平行场激励及 EMT 系统原型研制工作\upcite{xhl1998,tanyanting1999,lz2001}。随后，随着数字电路和可编程器件的发展，国内研究逐步实现了数字解调、实时采样和系统集成度提升\upcite{yin2010design,yin2011design,ywl2011electromagnetic}。
 
图 \ref{fig:ywl2010EMT_system} 所示为尹武良团队研制的数字 EMT 系统。该类研究推动了 EMT 系统由传统模拟平台向数字化、集成化方向发展。

\begin{figure}[htbp]
    \centering
    \includegraphics[width=0.58\textwidth]{figures/ywl2010.jpg}
    \figcaption{尹武良团队研制的数字 EMT 系统}{Digital EMT system developed by Professor Yin Wuliang's team}
    \label{fig:ywl2010EMT_system}
\end{figure}

 
在新型传感机制与硬件架构方面，国内研究也不断推进。例如，采用 TMR 传感器替代传统检测线圈的方案可提升局部磁场感知能力；基于 FPGA、STM32 和 ZYNQ 的系统架构则进一步增强了多通道采集、实时控制和数据处理能力\upcite{wangkai2018,zhengjiaoke2020,gaoyuan2024}。图 \ref{fig:gy2024EMT_system} 给出了高源团队研制的基于ZYNQ 的 EMT 系统。

\begin{figure}[htbp]
    \centering
    \includegraphics[width=0.58\textwidth]{figures/gy2024.jpg}
    \figcaption{高源团队研制的基于ZYNQ的EMT系统}{ZYNQ-based EMT system developed by Professor Gao Yuan's team}
    \label{fig:gy2024EMT_system}
\end{figure}

特别地，在轨道交通金属探伤领域，北京交通大学围绕钢轨探伤开展了系列专用化研究，包括面向强电磁干扰背景的钢轨探伤实验平台、适应钢轨轮廓的专用探头以及车载在线钢轨探伤系统等\upcite{HUO2021planner,liu2015electromagnetic,DZJY201611020,zhusheng2019}。这些研究在轨道场景 EMT 系统设计、微弱信号采集和成像算法验证方面积累了重要经验，也为本文面向道岔伤损检测的便携式系统设计提供了直接参考。
 
总体而言，国内外 EMT 系统研究已经在阵列传感器、激励采集硬件和专用检测平台方面取得较多进展，但面向高速铁路道岔这类空间受限、结构复杂且干扰强烈的特殊对象，仍缺乏兼顾便携性、高信噪比、多通道复数测量的完整解决方案。
 
\subsection{图像重建算法研究现状}
 
EMT 图像重建本质上是由有限边界测量数据反演内部介质参数分布的逆问题。由于 EMT 属于典型软场成像问题，场分布会随介质参数变化而改变，且边界测量维数通常远小于待重建空间维数，因此逆问题具有显著的病态性、非线性和噪声敏感性\upcite{WANG2020}。围绕这一问题，现有图像重建方法大体可分为传统灵敏度矩阵方法、滤波与状态估计方法以及数据驱动深度学习方法。
 
传统灵敏度矩阵方法通常基于小扰动线性近似，利用灵敏度矩阵建立边界测量变化与内部参数变化之间的线性关系。典型方法包括线性反投影（Linear Back Projection, LBP）、Tikhonov 正则化、Landweber 迭代、共轭梯度法和 Newton-Raphson 方法等。该类方法物理意义明确、实现简单，是 EMT 成像研究的重要基础；但其重建结果通常存在伪影明显、空间分辨率有限、对噪声敏感以及强非线性条件下精度下降等问题\upcite{mazhenqi2024,adler2011electrical}。
 
滤波与状态估计方法将 EMT 逆问题视为动态系统状态估计问题，通过卡尔曼滤波、扩展卡尔曼滤波或无迹卡尔曼滤波等方法融合多时刻观测信息，从而提高动态场景下的重建稳定性\upcite{zdy2025statistical,zhaoqian2022Kalman}。这类方法在时序信息利用和噪声抑制方面具有优势，但在高维、强非线性 EMT 问题中仍面临状态维数过高、模型构造复杂和参数调节困难等挑战。
 
随着深度学习的发展，数据驱动的非线性反演方法逐渐成为 EMT 图像重建的重要方向。深度学习方法可通过大量样本直接学习边界响应与内部参数分布之间的非线性映射，从而减少对线性近似和传统迭代求解的依赖。现有方法大致包括两类：一类是基于传统重建图像的后处理增强方法，例如利用 U-Net、GAN 等网络抑制伪影、提升分辨率；另一类是端到端反演方法，即直接由边界测量数据输出介质分布或缺陷图像\upcite{ronneberger2015unet,xiong2000electromagnetic,xiao2018deep,lixiaoyan2023improved,jiashuqing2025,carcavilla2025,fan2020solving}。后者更适合实时检测场景，也是本文后续模型设计的主要方向。
 
不过，纯数据驱动方法也存在一定局限：当训练数据覆盖不足或测试样本几何形状发生分布偏移时，模型容易学习到训练标签中的统计捷径，而未必真正满足底层物理规律。因此，如何将深度学习的非线性表达能力与 EMT 问题的电磁物理约束相结合，是当前智能反演方法进一步提升可靠性的关键。
 
\subsection{物理约束深度学习在电学成像中的研究进展}
 
近年来，物理约束深度学习逐渐成为电学成像和逆问题求解领域的重要研究方向。与纯数据驱动方法相比，这类方法试图在神经网络训练过程中引入控制方程、边界条件或中间物理变量，使模型在拟合观测数据的同时尽可能满足底层物理规律。Physics-Informed Neural Networks（PINN）是其中具有代表性的框架之一\upcite{raissi2019physics,karniadakis2021physics}。
 
在电阻抗成像（Electrical Impedance Tomography, EIT）等相关领域，已有研究表明，物理约束能够在一定程度上提升重建结果的稳定性、泛化能力和物理一致性。例如，Smyl 等人验证了在电学逆问题中引入物理方程对提高重建精度的积极作用；Yang 等人提出了 CNN 与 PINN 结合的两阶段重建框架；Zhu 等人则构建了包含电压预测与电导率重建分支的物理引导协同训练网络\upcite{smyl2025physics,yang2025two,zhu2026pgct}。
 
然而，在铁磁性 EMT 这类高导电率、高磁导率、强尺度差异和强非线性问题中，直接采用 PINN 主导反演面临显著挑战：一方面，材料参数跨度大与场分布的强非线性导致物理损失项量纲严重失衡，极易引发训练不稳定；另一方面，高阶导数计算成本高昂且对网格质量敏感。特别是在工程检测任务中，最终目标并非精确恢复完整电磁场，而是稳定检测伤损位置与范围。因此，相比完全依赖 PINN 框架，更合理的策略是以数据驱动反演为主线，在训练过程中引入物理场辅助监督或弱 PDE 正则化，以兼顾检测鲁棒性与物理一致性。
 
本文提出的 SCF-DeepONet 正是基于这一思路展开：以多视角复数电压响应作为输入，以二维伤损热图检测为主任务，逐步引入三维磁矢势差场辅助监督和 PDE 弱物理约束，通过这种``数据驱动为主线、物理信息渐进补强''的策略，SCF-DeepONet 在兼顾检测鲁棒性的同时，为 EMT 逆问题的可解释性与物理可靠性提供了层次化保障。

\section{本文主要研究内容与章节安排}

\subsection{本文主要研究内容}

针对高速铁路道岔表面及近表面伤损检测中存在的安装空间受限、微弱信号提取困难、传统 EMT 成像分辨率有限以及复杂工况下泛化能力不足等问题，本文围绕“面向道岔伤损检测任务的智能反演算法”与“便携式 EMT 硬件采集平台”两条主线展开研究，并通过软硬件集成实验构建完整的验证链路。主要内容如下。

（1）分析道岔伤损检测场景下 EMT 正问题与逆问题的理论基础。结合高导电率、高磁导率铁磁性目标特点，梳理传统基于灵敏度矩阵的成像方法及其局限性，重点讨论软场效应、逆问题病态性和线性化近似误差对道岔伤损检测的影响。

（2）提出面向道岔伤损检测的 SCF-DeepONet 智能反演模型，克服经典 DeepONet 在 EMT 多激励结构化输入和连续域输出上的适配不足。该模型以多视角复数差分电压响应为输入，以二维伤损热图为主输出，通过共享条件编码、Fourier 位置编码、FiLM 调制以及二维--三维分支耦合机制，在架构层面提升铁磁性 EMT 伤损检测的精度与鲁棒性。在此基础上，进一步引入三维磁矢势差场辅助任务，并通过多阶段训练策略实现二维检测与三维场表征的协同优化。

（3）围绕 SCF-DeepONet 开展模型优化与实验分析。通过消融实验验证所提出网络架构相对于经典 DeepONet 的有效性；构建纯热图监督与热图--场量联合监督两类基线，分析三维差场辅助任务对二维检测鲁棒性及分布外泛化能力的影响；进一步在三维场分支引入基于麦克斯韦方程相对残差的 PDE 弱物理约束，研究其对三维场重构物理一致性的改善效果。

（4）设计并实现一套便携式 EMT 硬件采集平台，解决微弱信号提取与多通道复数测量的工程难题。系统采用“测控板 + 信号切换板”的双板式结构，完成平面传感器阵列、收发异构多路选通网络、激励产生与功率放大、微弱信号调理与双相锁相解调、高精度模数转换、主控通信、电源管理和 PCB 抗干扰等模块设计，并在 100 kHz 主工作点下实现稳定的多通道复数互感响应采集。

（5）基于所搭建的硬件平台与上位机解析软件，开展系统级性能评估与仿损试块实验验证。通过激励稳定性、锁相放大链路、通道一致性等基础测试，以及传统成像方法与 SCF-DeepONet 的实测对比实验，验证便携式 EMT 系统从电磁采集到伤损成像的端到端可行性。
 
\subsection{本文章节安排}

全文共分为七章，各章内容安排如下。

第一章为绪论。主要介绍课题研究背景与工程意义，分析铁路道岔伤损检测的实际需求及现有主流探伤技术的局限性，综述 EMT 系统、图像重建算法和物理约束深度学习的国内外研究现状，并给出本文主要研究内容与章节安排。

第二章为电磁层析成像理论基础与传统成像方法。该章从 EMT 系统基本工作流程出发，阐述麦克斯韦方程基础上的正问题建模与有限元求解方法，分析 EMT 逆问题的数学表达及病态性来源，并介绍基于灵敏度矩阵的传统成像方法及其局限性。

第三章为基于 SCF-DeepONet 的道岔伤损电磁层析成像模型设计。该章结合道岔伤损检测任务特点，构建共享条件特征驱动的深度算子网络 SCF-DeepONet，介绍数据表示方式、网络总体架构、多任务损失函数以及四阶段训练策略。

第四章为 SCF-DeepONet 模型优化与实验结果分析。该章围绕模型训练与性能评估展开，使用经典DeepONet架构做消融实验，之后分析三维差场辅助任务对二维伤损检测和分布外泛化能力的影响，并研究 PDE 弱物理约束对三维场重构质量和物理一致性的改善效果。

第五章为便携式 EMT 系统硬件平台设计与实现。该章面向道岔伤损检测需求，介绍便携式 EMT 硬件平台的总体架构、平面线圈阵列、多路选通网络、激励与功率放大、锁相解调、高精度采集、通信模块、电源管理和 PCB 抗干扰设计。

第六章为系统性能评估与实验验证。该章基于所研制的便携式 EMT 系统，开展基础性能测试和仿损样本实验，测试了该系统在仿损样本上传统 EMT 成像方法和 SCF-DeepONet 模型的表现，验证了系统的可行性。

第七章为总结与展望。该章总结全文主要工作与创新点，分析研究中存在的不足与局限，并提出未来研究方向和改进建议。